\documentclass[11pt]{article}
\usepackage[margin=1in]{geometry}
\usepackage{amsmath,amssymb,bm}
\usepackage{booktabs}
\usepackage{array}
\usepackage{hyperref}
\usepackage{microtype}
\usepackage{siunitx}

\title{A One-Dimensional Correlated Configurational First-Passage Framework\\
for Fatigue Crack Initiation in Pure Aluminum}
\author{Donghyeok Kim \and Minsoo Kang \and Junhyeok Park}
\date{Theory Core v1 working manuscript, September 2026}

\newcommand{\Eop}{\mathrm{E}}

\begin{document}
\maketitle

\begin{abstract}
We formulate a one-dimensional fatigue crack-initiation framework in which microscopic interaction physics is retained as a constitutive foundation but the predicted observables remain macroscopic. A strict pure-normal, quasistatic spacing model cannot produce robust cycle-controlled accumulation without an additional slow state. We therefore introduce a correlated configurational coordinate $s_i$ coupled to local normal opening $a_i$ in an interacting two-row Lennard--Jones prototype. The full representative-region state is described by a many-body probability density $P_N(\bm a,\bm s,t)$, not by a product of independent one-cell distributions. Its overdamped dynamics is generated by the stress-dependent interaction energy and satisfies a nonnegative free-energy dissipation identity. Configurational motion is decomposed as $s_i=b n_i+\xi_i$, so well-index changes $n_i$ provide a natural record of irreversible rearrangement while $\xi_i$ denotes recoverable intrawell motion. Plastic and crack-opening instabilities are distinguished through the block Hessian: the relaxed configurational stiffness is the Schur complement $H_{ss}-H_{sa}H_{aa}^{-1}H_{as}$, whereas local opening stability is governed by $H_{aa}$. Fatigue crack initiation is defined as first passage through the opening saddle of the bound normal-energy basin rather than by an empirical scalar damage law. A dimensionless stochastic solver demonstrates the intended ordering with a single parameter set: low-load survival, configurational well crossing at intermediate load, opening-barrier reduction, and rapidly increasing first-passage probability at higher cyclic load. The calculation is a mechanism verification, not yet a calibrated aluminum fatigue-life prediction; characteristic length and an Al-specific EAM/MEAM energy calibration are intentionally deferred.
\end{abstract}

\section{Scope and scientific objective}
The objective is not to predict or measure an interatomic spacing distribution as an experimental end product. The microscopic state is hidden. The intended laboratory mapping is
\begin{equation}
 \sigma(0:t),T
 \longrightarrow
 \bar\varepsilon(t),\;
 \bar\varepsilon_p(t),\;
 \mathcal{D}(t),\;
 P_{\rm init}(t),
\end{equation}
where the predicted observables can be compared with cyclic tensile testing, stress--strain hysteresis, residual deformation, and crack-initiation time.

The previous strict pure-normal submodel establishes an important no-go result. If the structural normal coordinate follows the same reversible quasistatic branch every cycle, deterministic first passage saturates. If it instead undergoes fast stationary thermal renewal, integrated hazard per cycle scales primarily with elapsed time. Therefore cycle-controlled fatigue requires a physically derived slow microstructural state. The present Theory Core v1 identifies that state with a correlated configurational coordinate rather than an empirical scalar damage variable.

\section{Interacting two-row configurational geometry}
Consider $N$ local cells in a correlated representative region. The lower row is
\begin{equation}
 \bm r_j^-=(jb,0),
\end{equation}
and the upper-row cells are
\begin{equation}
 \bm r_i^+=\left(ib+\frac{b}{2}+s_i,\;a_i\right).
\end{equation}
Here $a_i$ is a local normal opening coordinate and $s_i$ is a hidden configurational coordinate. It is not introduced as a separately applied macroscopic shear strain. We decompose
\begin{equation}
 \boxed{s_i=b n_i+\xi_i},
 \qquad
 n_i\in\mathbb Z,\quad
 \xi_i\in[-b/2,b/2),
 \label{eq:sdecomp}
\end{equation}
where $n_i$ labels metastable registry wells and $\xi_i$ is intrawell motion.

The analytical prototype pair interaction is
\begin{equation}
 \phi_{\rm LJ}(r)
 =4\epsilon_{\rm LJ}
 \left[
 \left(\frac{\sigma_{\rm LJ}}{r}\right)^{12}
 -
 \left(\frac{\sigma_{\rm LJ}}{r}\right)^6
 \right].
\end{equation}
The many-body internal energy is constructed once from pair distances,
\begin{equation}
 \boxed{
 \mathcal U_N(\bm a,\bm s)
 =
 \sum_{\alpha<\beta}
 \phi_{\rm LJ}\!\left(R_{\alpha\beta}(\bm a,\bm s)\right).
 }
 \label{eq:UN}
\end{equation}
This avoids double counting a separate normal-chain energy and a two-row energy. Previously derived normal-chain and registry surfaces are interpreted as restrictions or projections of Eq.~\eqref{eq:UN}, not additive independent energies.

\section{External normal work and macroscopic strain bridge}
For a selected configurational mode, define the representative-region axial strain map
\begin{equation}
 \varepsilon(\bm a,\bm s)
 =
 \frac{1}{Na_0}
 \sum_{i=1}^{N}
 \left[(a_i-a_0)+\chi s_i\right].
 \label{eq:strainmap}
\end{equation}
The coefficient $\chi$ is an explicit geometry/coarse-graining projection. It is not a fatigue-damage fitting coefficient. In a strict registry mode with no axial projection one sets $\chi=0$; in an inclined or otherwise axially projecting configurational mode it must be derived or calibrated from that geometry.

With representative volume $V_c$, the stress-controlled configurational potential is
\begin{equation}
 \boxed{
 \mathcal G_N(\bm a,\bm s,t)
 =
 \mathcal U_N(\bm a,\bm s)
 -
 V_c\sigma(t)\varepsilon(\bm a,\bm s).
 }
 \label{eq:GN}
\end{equation}
The macroscopic strain is the ensemble average
\begin{equation}
 \boxed{
 \bar\varepsilon(t)
 =
 \int \varepsilon(\bm a,\bm s)
 P_N(\bm a,\bm s,t)\,d\bm a\,d\bm s.
 }
 \label{eq:macrostrain}
\end{equation}

\section{Correlated stochastic dynamics}
Let
\begin{equation}
 \bm q=(a_1,\ldots,a_N,s_1,\ldots,s_N).
\end{equation}
The foundational probability density is
\begin{equation}
 \boxed{P_N(\bm q,t)},
\end{equation}
with no product closure
\begin{equation}
 P_N(\bm q,t)\neq\prod_i P_i(a_i,s_i,t)
\end{equation}
assumed at the theory level.

For diagonal positive mobility matrix
\begin{equation}
 \bm M=\mathrm{diag}(M_a,\ldots,M_a,M_s,\ldots,M_s),
\end{equation}
the coarse-grained overdamped configurational dynamics is
\begin{equation}
 d\bm q
 =
 -\bm M\nabla_{\bm q}\mathcal G_N\,dt
 +
 \sqrt{2k_BT\bm M}\,d\bm W_t .
 \label{eq:langevin}
\end{equation}
Its many-body Smoluchowski equation is
\begin{equation}
 \boxed{
 \frac{\partial P_N}{\partial t}
 =
 \nabla_{\bm q}\cdot
 \left[
 \bm M
 \left(
 P_N\nabla_{\bm q}\mathcal G_N
 +k_BT\nabla_{\bm q}P_N
 \right)
 \right].
 }
 \label{eq:smol}
\end{equation}
Equation~\eqref{eq:smol}, rather than an Arrhenius jump law, is the governing dynamics. An Arrhenius or Kramers rate is only a possible high-barrier reduced description after a metastable basin and saddle have been derived from the same energy landscape.

Define
\begin{equation}
 \mathcal A[P_N]
 =
 \int P_N\mathcal G_N\,d\bm q
 +
 k_BT\int P_N\ln P_N\,d\bm q.
\end{equation}
With probability flux
\begin{equation}
 \bm J
 =
 -\bm M\left(P_N\nabla_{\bm q}\mathcal G_N+k_BT\nabla_{\bm q}P_N\right),
\end{equation}
the irreversible dissipation is
\begin{equation}
 \boxed{
 \mathcal D
 =
 \int \frac{\bm J^T\bm M^{-1}\bm J}{P_N}\,d\bm q\ge0.
 }
 \label{eq:dissipation}
\end{equation}

\section{Plastic memory and hysteresis}
Equation~\eqref{eq:sdecomp} separates reversible and irreversible configurational motion. A well crossing changes $n_i$, whereas intrawell relaxation changes $\xi_i$. A convenient coarse-grained plastic contribution associated with the selected axial projection is
\begin{equation}
 \bar\varepsilon_p(t)
 =
 \frac{\chi b}{Na_0}\sum_i \Eop[n_i(t)].
 \label{eq:plastic}
\end{equation}
More generally, when the selected configurational geometry has no direct axial projection, residual plastic strain is computed by unloading the current configurational state and minimizing the normal coordinates; the resulting zero-stress relaxed length provides the residual strain.

Because $M_s$ is finite and generally smaller than the fast normal mobility, the configurational distribution need not return to its initial state after a closed stress cycle:
\begin{equation}
 P_N(t+T_f)\neq P_N(t).
\end{equation}
Consequently loading and unloading can follow different macroscopic paths,
\begin{equation}
 \oint \sigma\,d\bar\varepsilon\neq0,
\end{equation}
without imposing a phenomenological hysteresis rule.

\section{Separate plastic and opening instabilities}
Let the Hessian of $\mathcal G_N$ be partitioned as
\begin{equation}
 \bm H=
 \begin{pmatrix}
 H_{aa} & H_{as}\\
 H_{sa} & H_{ss}
 \end{pmatrix}.
\end{equation}
If $a$ relaxes rapidly while $s$ evolves slowly, the relaxed configurational stiffness is the Schur complement
\begin{equation}
 \boxed{
 H_{\rm pl}=H_{ss}-H_{sa}H_{aa}^{-1}H_{as}.
 }
 \label{eq:Hpl}
\end{equation}
An athermal plastic spinodal is therefore identified by
\begin{equation}
 \lambda_{\min}(H_{\rm pl})=0.
\end{equation}

Crack opening is a distinct instability. On the fast opening time scale the configurational coordinate is frozen, so the instantaneous opening stiffness is governed by
\begin{equation}
 \boxed{H_{\rm open}=H_{aa}}.
\end{equation}
The athermal opening limit satisfies
\begin{equation}
 \lambda_{\min}(H_{aa})=0.
\end{equation}
This modal distinction prevents a configurational soft mode from being misclassified as a crack-opening instability.

\section{Opening saddle and first-passage crack initiation}
For a local opening coordinate $a$, frozen configurational state $s$, and local tensile force $F$, define
\begin{equation}
 g(a;s,F)=U_{\rm loc}(a,s)-F(a-a_0).
\end{equation}
The intact minimum $a_m$ and opening saddle $a^\ddagger$ satisfy
\begin{align}
 \partial_a g(a_m;s,F)&=0,
 &\partial_a^2g(a_m;s,F)&>0,\\
 \partial_a g(a^\ddagger;s,F)&=0,
 &\partial_a^2g(a^\ddagger;s,F)&<0.
\end{align}
The mechanically derived opening barrier is
\begin{equation}
 \boxed{
 \Delta G_c^\ddagger(s,F)
 =g(a^\ddagger;s,F)-g(a_m;s,F).
 }
 \label{eq:barrier}
\end{equation}
The dividing surface is the opening saddle, not an arbitrary spacing threshold.

Let $\Omega_b(t)$ denote the bound basin and $\Gamma_c(t)$ its opening-dividing surface. Local survival and crack-initiation probability are
\begin{equation}
 S_c(t)=\int_{\Omega_b(t)}P_N(\bm q,t)\,d\bm q,
 \qquad
 \boxed{P_{\rm init}^{(c)}(t)=1-S_c(t).}
 \label{eq:survival}
\end{equation}
For a fixed boundary the escape rate is the probability flux
\begin{equation}
 -\dot S_c=\int_{\Gamma_c}\bm J\cdot\bm n\,dS.
 \label{eq:flux}
\end{equation}
Thus fatigue history is stored in the evolving nonequilibrium distribution, not in a separately postulated scalar damage variable.

The characteristic length or area required to aggregate correlated local regions into a specimen-scale weakest-link probability is intentionally deferred. Theory Core v1 ends at the correlated representative-region level.

\section{Executable solver}
The accompanying \texttt{solver\_v1} package integrates Eq.~\eqref{eq:langevin} directly with Euler--Maruyama trajectories. Correlation is retained by integrating the full interacting state; no one-cell product closure is used. The local opening saddle is obtained by precomputing the normal traction curve $\partial U_{\rm loc}/\partial a$ over registry state and force, then locating its stable and unstable branches. A trajectory is absorbed when a cell crosses the mechanically derived opening saddle or loses the bound minimum.

The solver returns the macroscopic strain, stress--strain loop, configurational well-crossing activity, opening barrier, survival, and first-passage fraction.

\section{Dimensionless mechanism verification}
A proof-of-principle screening was performed with one fixed dimensionless parameter set:
\begin{equation}
 N=3,\quad b=1,\quad
 \epsilon_{\rm LJ}=1,\quad
 \sigma_{\rm LJ}=0.82,\quad
 k_BT=0.009,
\end{equation}
\begin{equation}
 M_a=1,\quad M_s=0.15,\quad \chi=0.40,
\end{equation}
using 32 stochastic trajectories, 10 cycles, and time step $\Delta t=0.02$. These values are not claimed as aluminum material parameters.

\begin{center}
\begin{tabular}{cccc}
\toprule
$F_{\max}$ & first-passage fraction & max. mean $|n_i|$ & min. mean $\Delta G_c^\ddagger$ \\
\midrule
2.5 & 0 & 0 & 0.3549 \\
3.2 & 0 & 0.3021 & 0.1062 \\
3.4 & 0.09375 & 3.1839 & 0.0609 \\
3.6 & 0.84375 & 6.8000 & 0.0271 \\
\bottomrule
\end{tabular}
\end{center}

The key result is the ordering rather than the numerical values. At $F_{\max}=2.5$, the model remains effectively within the initial configurational well and no first passage occurs. At $F_{\max}=3.2$, configurational well crossing is already present while the opening basin survives all trajectories in the finite ensemble. At $F_{\max}=3.4$, the mean opening barrier is reduced and a nonzero first-passage population appears. At $F_{\max}=3.6$, opening first passage becomes dominant. The same solver also produces nonzero cyclic stress--strain loop area. At the highest load the survivor-conditioned loop area is not expected to remain monotonic because most trajectories have already been absorbed.

This numerical experiment therefore verifies that the proposed architecture can produce the desired sequence
\begin{equation}
 \boxed{
 \text{cyclic loading}
 \rightarrow
 \text{configurational rearrangement}
 \rightarrow
 \text{hysteresis/plastic memory}
 \rightarrow
 \Delta G_c^\ddagger\downarrow
 \rightarrow
 P_{\rm init}\uparrow
 }.
\end{equation}

\section{Relation to the pure-normal no-go result}
The slow configurational state resolves the structural limitation of the strict normal-spacing model. The system is no longer required to revisit the same internal distribution after each identical macroscopic stress cycle. Cycle history modifies $P_N(\bm a,\bm s,t)$, which in turn modifies the opening barrier and the next-cycle response. Therefore cycle accumulation is generated by internal-state evolution rather than by assigning an empirical cycle counter to an otherwise reversible normal coordinate.

This does not establish that the present LJ configurational coordinate is the unique or quantitatively correct aluminum plastic mechanism. Published room-temperature aluminum studies instead emphasize persistent/secondary slip, dislocation structure, and orientation-dependent initiation. Those observations motivate the present slow-state extension and also require future replacement or calibration of the analytical LJ prototype by an Al-appropriate energy landscape.

\section{Limitations and next calibration steps}
Theory Core v1 deliberately separates structural choices from parameters that remain unresolved.
\begin{enumerate}
 \item The LJ interaction is an analytical prototype. Quantitative aluminum prediction requires EAM/MEAM or validated atomistic energy/barrier data.
 \item The geometry coefficient $\chi$ must be derived from the selected configurational mode or independently calibrated; it must not be tuned to an S--N curve.
 \item Mobilities $M_a$ and $M_s$ require time-scale calibration.
 \item The present local opening saddle freezes the configurational coordinate. A later calculation may replace this by a full multidimensional minimum-energy-path dividing surface.
 \item Characteristic correlation length/area and specimen-scale aggregation are deferred.
 \item The reported stochastic screening demonstrates mechanism viability, not pure-Al lifetime accuracy.
\end{enumerate}

\section{Conclusion}
The final Theory Core v1 is a correlated, one-dimensional normal-loading fatigue-initiation framework with a hidden configurational plastic state. Its foundational equation is the many-body Smoluchowski dynamics generated by the same interaction energy used to define stability and crack-opening barriers. Plastic and crack modes are separated by Hessian structure, and crack initiation is a first-passage event through a derived opening saddle. The model does not require independent-atom probability closure, an imposed hysteresis rule, or an empirical scalar damage law.

The dimensionless solver verifies the central mechanism with one parameter set: configurational well crossing can precede crack initiation, cyclic rearrangement reduces the opening barrier, and the first-passage fraction increases strongly with load. The theory is therefore sufficiently closed to serve as the project's fixed v1 backbone. Subsequent work should refine material calibration and spatial correlation without changing this core state--dynamics--stability--first-passage architecture unless quantitative validation falsifies it.

\begin{thebibliography}{9}
\bibitem{DanielsDorn1957}
N. H. G. Daniels and J. E. Dorn,
``The Effect of Temperature, Frequency, and Grain Size on the Fatigue Properties of High-Purity Aluminum,''
ASTM STP 196, p. 94 (1957), DOI: 10.1520/STP19619570007.

\bibitem{Zhai1996}
T. Zhai, G. A. D. Briggs, and J. W. Martin,
``Fatigue damage at room temperature in aluminium single crystals---IV. Secondary slip,''
\emph{Acta Materialia} 44(9), 3489--3496 (1996), DOI: 10.1016/1359-6454(96)00025-0.

\bibitem{Hayashi2022}
M. Hayashi,
``Effect of crystal orientation on fatigue crack initiation life in pure aluminum single crystals,''
\emph{International Journal of Fatigue} 156, 106661 (2022), DOI: 10.1016/j.ijfatigue.2021.106661.
\end{thebibliography}

\end{document}
